\documentclass[12pt]{article}
\usepackage[utf8]{inputenc}
\usepackage[T1]{fontenc}
\usepackage{hyperref}
\title{Udmurt Language Resources for K-12 Curriculum Development}
\author{Compiled by Claude Sonnet 5 (claude-sonnet-5) for Andr\'as Kornai}
\date{}
\begin{document}
\maketitle

\section{Introduction}
Udmurt (ISO 639-3 udm; Glottocode udmu1245) is a Uralic language of the Russian Federation. This report is based on a small number of web searches, so it is short and incomplete. Items marked as not verified were mentioned in search results but not opened.

\section{Primary Educational Resources}
\begin{itemize}
\item The school textbook series ``Z\v{e}ch-a, bur-a, udmurt kyl!'' (written in Cyrillic in the original) for grades 7--9 is compared with a Swedish series in a journal article on electronic textbooks. I found no publisher or ISBN details. \url{https://journals.rudn.ru/polylinguality/article/view/32868/en_US}
\item An article from 2019 on the possibilities of teaching Udmurt in the Udmurt Republic covers teaching aids and electronic textbooks. \url{https://gramota.net/eng/materials/4/2019/2/6.html}
\item The textbook \emph{Marym, l\'esya} (2005) takes a communicative approach and is suited to courses taught by native speakers. LMU Munich's INFUSE project adapts it as an open e-learning course with English grammar explanations. \url{https://www.infuse.finnougristik.uni-muenchen.de/e-learning/udmurt/index.html}
\end{itemize}

\section{Linguistic and Reference Materials}
\begin{itemize}
\item Freelang's free online Udmurt--English and English--Udmurt dictionary is small, about 1,500 words, and was last updated in 2014. \url{https://www.freelang.net/dictionary/udmurt.php}
\item A paper describes a corpus-oriented dictionary database for Beserman, a dialect of Udmurt. \url{https://real.mtak.hu/65045/}
\item The Volga-Kama corpora include a main Udmurt corpus and a spoken Udmurt corpus. A paper gives the address below, and I did not verify that the site is online. \url{http://volgakama.web-corpora.net/}
\item A type design article describes the Udmurt Cyrillic alphabet and its extended letters. \url{https://type.today/en/journal/udmurt}
\end{itemize}

\section{Religious and Community Resources}
My searches found nothing on Udmurt scripture or community organisations, so I list none.

\section{Digital and Audio Resources}
\begin{itemize}
\item A 2018 paper describes a sound-aligned corpus of dialectal Udmurt speech, about 3.5 hours at the time, with a web interface. \url{https://preview.aclanthology.org/moar-dois/W18-0203.pdf}
\end{itemize}

\section{Curriculum Development Context}
Udmurt-language schoolbooks were produced in the 1920s and 1930s, and production stopped from the 1960s. A later national-regional curriculum component gave schools room to reintroduce the language. Mercator's regional dossier describes this history. \url{https://wiki.mercator-research.eu/languages:udmurt_in_russia}
The searches did not turn up a full native-language K-12 curriculum. Textbooks for grades 1--6 and 10--11, and any current policy, need further searching.

\end{document}